\section{Maximum rank: one transposition for each tree edge}
\begin{theorem}\label{thm:maxrank}
Let $n\geq7$ and let $G\leq S_n$ have a C-group marking of rank $n-1$.
Then $G=S_n$, and its distinguished generators are exactly the
transpositions on the edges of a tree. Conversely every such tree
gives a C-group marking. Its Coxeter diagram is the line graph of the
tree, and its coset system is a regular hypertope with full
type-preserving group $S_n$.
\end{theorem}
\begin{proof}
Apply UB.Extremal to the independent family of $n-1$ elements.
It is either an ordinary tree marking or the sign marking of UB.Sign.
In the latter case, an edge $e$ adjacent to its distinguished edge $d$
exists because the tree has at least three vertices and is connected.
The mark $(\tau_d\tau_e)^{\pm1}$ has order three, contradicting the
involution hypothesis. Thus the marking is ordinary.

UB.F gives its complete intersections and flags, its trivial Borel
and its panels of size two. Edge transpositions have product order
three when their edges meet and order two otherwise, giving the line
graph. The common coset criterion gives residual connectedness.
The group acts regularly on chambers; their graph is connected because
the marks generate. A type-preserving automorphism fixing a chamber
fixes its unique adjacent chamber of each type, hence every chamber
and every object. The full type-preserving group is therefore $S_n$.
\end{proof}
The common tree-presentation theorem now recovers [FL, Proposition 3.1],
with the line-graph Coxeter relations and the triangle relations
$(s_es_fs_es_g)^2=1$. Its proof holds for every finite tree, even below
$n=7$; the classification retains its stated degree range.
